%%%%%%%%%%%%%%%%%%%%%%%%%%%%%%%%%%%%%%%%%%%%%%%%%%%%%%%%%%%%%%%
%
% Welcome to Overleaf --- just edit your LaTeX on the left,
% and we'll compile it for you on the right. If you open the
% 'Share' menu, you can invite other users to edit at the same
% time. See www.overleaf.com/learn for more info. Enjoy!
%
%%%%%%%%%%%%%%%%%%%%%%%%%%%%%%%%%%%%%%%%%%%%%%%%%%%%%%%%%%%%%%%
\documentclass{article}
\usepackage{geometry}
\usepackage{url}
\geometry{landscape, textwidth=27cm}
\usepackage{circuitikz}
\begin{document}

\newcommand{\bipole}[1]{%
\texttt{#1}\kern5pt \hfill
\begin{circuitikz}%
\draw (0,0) to[ #1 ] (2,0); 
\end{circuitikz}%
\hspace{5mm}%
}

\renewcommand{\arraystretch}{3}
\section{Overall Examples}
\begin{center}
\begin{tabular}{lll} 
\bipole{empty diode} & \bipole{empty Schottky diode} & \bipole{empty Zener diode}\\
\bipole{empty tunnel diode} & \bipole{photodiode} & \bipole{empty diode}\\
\bipole{empty varcap} & \bipole{full diode} & \bipole{full Schottky diode}\\
\bipole{full Zener diode} & \bipole{full tunnel diode} & \bipole{full photodiode}\\

\bipole{barrier} & & \\
\end{tabular}
\end{center}


\section{Electrical Circuit}    
\begin{center}
\begin{circuitikz}[american voltages]
\draw
  (0,0) to [short, *-] (6,0)
  to [V, l_=$\mathrm{j}{\omega}_m \underline{\psi}^s_R$] (6,2) 
  to [R, l_=$R_R$] (6,4) 
  to [short, i_=$\underline{i}^s_R$] (5,4) 
  (0,0) to [open, v^>=$\underline{u}^s_s$] (0,4) 
  to [short, *- ,i=$\underline{i}^s_s$] (1,4) 
  to [R, l=$R_s$] (3,4)
  to [L, l=$L_{\sigma}$] (5,4) 
  to [short, i_=$\underline{i}^s_M$] (5,3) 
  to [L, l_=$L_M$] (5,0); 
\end{circuitikz}
\end{center}
% \end{document}\\
For additional examples with TikZ, visit \url{https://tikz.net/}, whereas particularly for electronics, explore the following: \url{https://tikz.net/electric_circuit_resistor/} 
\end{document}